% TT30002C
% Verbrennung, schwer
% 14.0
% 2013-11-30
\label{m:verbrennung-schwer}
\begin{itemize}
  \item \TxMassVitMK
  \item Kleidung und Schmuck entfernen (sofern nicht  eingeschmolzen!)
  \item Besonders auf \EE{Wärmeerhalt} achten!
  \item Luftzug vermeiden
  \item \Z{\E{Keine} Kühlung\footnote{Bei großflächigen
  Verbrennungen sollte der Patient nur primär abgelöscht werden. Hier soll
  keine Kühlung der betroffenen Hautregionen erfolgen.
  \cite{Asanger2010}}}
%   \item Mind. 10 min. mit mäßig kaltem Wasser die betroffenen Körperregionen spülen  --
%   \EE{NICHT unterkühlen!} 
%   OBSOLETE: Großflächige Verbrennungen (Erwachsene
% > 20% KOF, Kinder > 10% KOF, Säuglinge > 5%) sollten nur primär
% abgelöscht und keimfrei abgedeckt werden. Hier erfolgt keine Kühlung
% der betroffenen Regionen.
  \item \EE{Abdeckung}: Sterile, nicht verklebende
  Wundauflagen, evtl. metalliserte Wundauflagen
  \cite{Helfen2008,RedelsteinerHRNS1}.\ACh{GABS}{2010-08-06}{Metallisierte
  Verbände sind wünschenswert, werden zB in Helfen2008 
  empfohlen.}\ZFootnote{Der Nutzen  von kühlenden
  Spezialverbänden, z.\,B. Waterjel{\texttrademark} ist
  aufgrund  der Unterkühlungs-Gefahr und einer  möglichen
  Durchblutungsstörung der Wunde
  (Gefahr von Wundheilungsstörungen im weiteren Verlauf)
  umstritten.} Locker anlegen. Bei großflächigen
  Verbrennungen können,  wenn das sterile Material nicht
  ausreicht,  zur Not auch saubere Leintücher o.\,ä. 
  verwendet werden \cite{emergency9,Asanger2010}.
  % TODO INHALT: Stellungsnahme zu Water-Jel
  %         \item Wenn vorhanden: Water-Jel
  %         \footnote{Streitfrage: Water-Jel: Umstritten}
  \item Schockbekämpfung
  \item \Z{NKV+: Infusionstherapie mit \E{kristalloiden} Lösungen; kolloidale
  Lösungen sind \E{kontraindiziert}.}
\end{itemize}